% Options for packages loaded elsewhere
\PassOptionsToPackage{unicode}{hyperref}
\PassOptionsToPackage{hyphens}{url}
\documentclass[
]{article}
\usepackage{xcolor}
\usepackage[margin=1in]{geometry}
\usepackage{amsmath,amssymb}
\setcounter{secnumdepth}{-\maxdimen} % remove section numbering
\usepackage{iftex}
\ifPDFTeX
  \usepackage[T1]{fontenc}
  \usepackage[utf8]{inputenc}
  \usepackage{textcomp} % provide euro and other symbols
\else % if luatex or xetex
  \usepackage{unicode-math} % this also loads fontspec
  \defaultfontfeatures{Scale=MatchLowercase}
  \defaultfontfeatures[\rmfamily]{Ligatures=TeX,Scale=1}
\fi
\usepackage{lmodern}
\ifPDFTeX\else
  % xetex/luatex font selection
\fi
% Use upquote if available, for straight quotes in verbatim environments
\IfFileExists{upquote.sty}{\usepackage{upquote}}{}
\IfFileExists{microtype.sty}{% use microtype if available
  \usepackage[]{microtype}
  \UseMicrotypeSet[protrusion]{basicmath} % disable protrusion for tt fonts
}{}
\makeatletter
\@ifundefined{KOMAClassName}{% if non-KOMA class
  \IfFileExists{parskip.sty}{%
    \usepackage{parskip}
  }{% else
    \setlength{\parindent}{0pt}
    \setlength{\parskip}{6pt plus 2pt minus 1pt}}
}{% if KOMA class
  \KOMAoptions{parskip=half}}
\makeatother
\usepackage{color}
\usepackage{fancyvrb}
\newcommand{\VerbBar}{|}
\newcommand{\VERB}{\Verb[commandchars=\\\{\}]}
\DefineVerbatimEnvironment{Highlighting}{Verbatim}{commandchars=\\\{\}}
% Add ',fontsize=\small' for more characters per line
\usepackage{framed}
\definecolor{shadecolor}{RGB}{248,248,248}
\newenvironment{Shaded}{\begin{snugshade}}{\end{snugshade}}
\newcommand{\AlertTok}[1]{\textcolor[rgb]{0.94,0.16,0.16}{#1}}
\newcommand{\AnnotationTok}[1]{\textcolor[rgb]{0.56,0.35,0.01}{\textbf{\textit{#1}}}}
\newcommand{\AttributeTok}[1]{\textcolor[rgb]{0.13,0.29,0.53}{#1}}
\newcommand{\BaseNTok}[1]{\textcolor[rgb]{0.00,0.00,0.81}{#1}}
\newcommand{\BuiltInTok}[1]{#1}
\newcommand{\CharTok}[1]{\textcolor[rgb]{0.31,0.60,0.02}{#1}}
\newcommand{\CommentTok}[1]{\textcolor[rgb]{0.56,0.35,0.01}{\textit{#1}}}
\newcommand{\CommentVarTok}[1]{\textcolor[rgb]{0.56,0.35,0.01}{\textbf{\textit{#1}}}}
\newcommand{\ConstantTok}[1]{\textcolor[rgb]{0.56,0.35,0.01}{#1}}
\newcommand{\ControlFlowTok}[1]{\textcolor[rgb]{0.13,0.29,0.53}{\textbf{#1}}}
\newcommand{\DataTypeTok}[1]{\textcolor[rgb]{0.13,0.29,0.53}{#1}}
\newcommand{\DecValTok}[1]{\textcolor[rgb]{0.00,0.00,0.81}{#1}}
\newcommand{\DocumentationTok}[1]{\textcolor[rgb]{0.56,0.35,0.01}{\textbf{\textit{#1}}}}
\newcommand{\ErrorTok}[1]{\textcolor[rgb]{0.64,0.00,0.00}{\textbf{#1}}}
\newcommand{\ExtensionTok}[1]{#1}
\newcommand{\FloatTok}[1]{\textcolor[rgb]{0.00,0.00,0.81}{#1}}
\newcommand{\FunctionTok}[1]{\textcolor[rgb]{0.13,0.29,0.53}{\textbf{#1}}}
\newcommand{\ImportTok}[1]{#1}
\newcommand{\InformationTok}[1]{\textcolor[rgb]{0.56,0.35,0.01}{\textbf{\textit{#1}}}}
\newcommand{\KeywordTok}[1]{\textcolor[rgb]{0.13,0.29,0.53}{\textbf{#1}}}
\newcommand{\NormalTok}[1]{#1}
\newcommand{\OperatorTok}[1]{\textcolor[rgb]{0.81,0.36,0.00}{\textbf{#1}}}
\newcommand{\OtherTok}[1]{\textcolor[rgb]{0.56,0.35,0.01}{#1}}
\newcommand{\PreprocessorTok}[1]{\textcolor[rgb]{0.56,0.35,0.01}{\textit{#1}}}
\newcommand{\RegionMarkerTok}[1]{#1}
\newcommand{\SpecialCharTok}[1]{\textcolor[rgb]{0.81,0.36,0.00}{\textbf{#1}}}
\newcommand{\SpecialStringTok}[1]{\textcolor[rgb]{0.31,0.60,0.02}{#1}}
\newcommand{\StringTok}[1]{\textcolor[rgb]{0.31,0.60,0.02}{#1}}
\newcommand{\VariableTok}[1]{\textcolor[rgb]{0.00,0.00,0.00}{#1}}
\newcommand{\VerbatimStringTok}[1]{\textcolor[rgb]{0.31,0.60,0.02}{#1}}
\newcommand{\WarningTok}[1]{\textcolor[rgb]{0.56,0.35,0.01}{\textbf{\textit{#1}}}}
\usepackage{longtable,booktabs,array}
\newcounter{none} % for unnumbered tables
\usepackage{calc} % for calculating minipage widths
% Correct order of tables after \paragraph or \subparagraph
\usepackage{etoolbox}
\makeatletter
\patchcmd\longtable{\par}{\if@noskipsec\mbox{}\fi\par}{}{}
\makeatother
% Allow footnotes in longtable head/foot
\IfFileExists{footnotehyper.sty}{\usepackage{footnotehyper}}{\usepackage{footnote}}
\makesavenoteenv{longtable}
\usepackage{graphicx}
\makeatletter
\newsavebox\pandoc@box
\newcommand*\pandocbounded[1]{% scales image to fit in text height/width
  \sbox\pandoc@box{#1}%
  \Gscale@div\@tempa{\textheight}{\dimexpr\ht\pandoc@box+\dp\pandoc@box\relax}%
  \Gscale@div\@tempb{\linewidth}{\wd\pandoc@box}%
  \ifdim\@tempb\p@<\@tempa\p@\let\@tempa\@tempb\fi% select the smaller of both
  \ifdim\@tempa\p@<\p@\scalebox{\@tempa}{\usebox\pandoc@box}%
  \else\usebox{\pandoc@box}%
  \fi%
}
% Set default figure placement to htbp
\def\fps@figure{htbp}
\makeatother
\setlength{\emergencystretch}{3em} % prevent overfull lines
\providecommand{\tightlist}{%
  \setlength{\itemsep}{0pt}\setlength{\parskip}{0pt}}
\usepackage{bookmark}
\IfFileExists{xurl.sty}{\usepackage{xurl}}{} % add URL line breaks if available
\urlstyle{same}
% fallback for those not using the hyperref driver hyperxmp:
\makeatletter
\@ifundefined{xmpquote}{\newcommand{\xmpquote}[1]{#1}}{}
\makeatother
\hypersetup{
  pdftitle={Workflows and Write Ups},
  pdfauthor={Amanda Mae Woodward},
  hidelinks,
  pdfcreator={LaTeX via pandoc}}

\title{Workflows and Write Ups}
\author{Amanda Mae Woodward}
\date{2026-11-17}

\begin{document}
\maketitle

The goal of this class is to help us figure out how to present our
results.

\subsection{Learning Outcomes:}\label{learning-outcomes}

By the end of this class, students should be able to:

\begin{itemize}
\item
  Describe modifications that can be made to R code chunks
\item
  Explain how to write code inline with text
\item
  Create Write Ups that combine graphs, interpretation, and analyses
  into a streamlined document.
\end{itemize}

\subsection{Learning Outcome \#1: Describe modifications that can be
made to R code
chunks}\label{learning-outcome-1-describe-modifications-that-can-be-made-to-r-code-chunks}

Today we're going to focus on code chunks themselves, rather than what
goes into a code chunk. You'll notice that I will be switching between
knitted versions and non knitted versions throughout class. This is to
help us get a sense of what these options do to our codes.

\textbf{Disclaimer:} This is going to be a brief introduction. You can
write full APA style papers in R Markdown files, and this is typically
easier if you use LaTex code. If you are interested in learning more
about this option, here is a good resource for it: psychstatsworkshop

\subsubsection{Changing Code Chunks}\label{changing-code-chunks}

When we are making documents in R Markdown, we don't always want people
to see all of our code and all of our output. We can modify our code
chunks to make ``hide'' certain aspects of code.

\begin{enumerate}
\def\labelenumi{\arabic{enumi}.}
\tightlist
\item
  \textbf{echo:}
\end{enumerate}

\begin{itemize}
\tightlist
\item
  if TRUE, shows code in knitted file
\item
  if FALSE, does not show code
\end{itemize}

\begin{Shaded}
\begin{Highlighting}[]
\FunctionTok{data}\NormalTok{(mtcars)}
\FunctionTok{mean}\NormalTok{(mtcars}\SpecialCharTok{$}\NormalTok{mpg)}
\end{Highlighting}
\end{Shaded}

\begin{verbatim}
## [1] 20.09062
\end{verbatim}

\begin{verbatim}
## [1] 20.09062
\end{verbatim}

\begin{enumerate}
\def\labelenumi{\arabic{enumi}.}
\setcounter{enumi}{1}
\tightlist
\item
  \textbf{include:}
\end{enumerate}

\begin{itemize}
\tightlist
\item
  if TRUE, code and output are included in the file
\item
  if FALSE, code and output are not included in the file
\end{itemize}

\begin{Shaded}
\begin{Highlighting}[]
\FunctionTok{library}\NormalTok{(ggplot2)}
\end{Highlighting}
\end{Shaded}

\begin{verbatim}
## Warning: package 'ggplot2' was built under R version 4.5.2
\end{verbatim}

\begin{Shaded}
\begin{Highlighting}[]
\FunctionTok{ggplot}\NormalTok{(mtcars, }\FunctionTok{aes}\NormalTok{(mpg, hp))}\SpecialCharTok{+}\FunctionTok{geom\_point}\NormalTok{()}\SpecialCharTok{+}\FunctionTok{theme\_classic}\NormalTok{()}
\end{Highlighting}
\end{Shaded}

\pandocbounded{\includegraphics[keepaspectratio]{workflows-and-writeups_files/figure-latex/testIncludeTRUE-1.pdf}}

\begin{enumerate}
\def\labelenumi{\arabic{enumi}.}
\setcounter{enumi}{2}
\tightlist
\item
  \textbf{error:}
\end{enumerate}

\begin{itemize}
\tightlist
\item
  if TRUE, prints error messages in document
\item
  if FALSE, does not print error messages
\end{itemize}

\begin{enumerate}
\def\labelenumi{\arabic{enumi}.}
\setcounter{enumi}{3}
\tightlist
\item
  \textbf{message and warning:}
\end{enumerate}

\begin{itemize}
\tightlist
\item
  two separate options
\item
  if TRUE, includes R messages and warnings in your document
\item
  if FALSE, does not include R messages or warnings
\end{itemize}

\subsubsection{Practice:}\label{practice}

Try these on your own! Create a code chunk with a graph of your favorite
data this semester.

\begin{enumerate}
\def\labelenumi{\arabic{enumi}.}
\item
  What happens if you use echo= FALSE?
\item
  What happens if you use include FALSE?
\item
  Try combining multiple arguments in your code chunk (pick your
  favorites).
\end{enumerate}

\subsubsection{More resources:}\label{more-resources}

The R/Posit Cloud Cheatsheets have a whole list of things you can
include or not.

\subsection{Learning Outcome \#2: Explain how to write code inline with
text}\label{learning-outcome-2-explain-how-to-write-code-inline-with-text}

You can also use r code inline (meaning in your actual sentences),
rather than in code chunks. The way to set this up is to use \texttt{}.

If you want to highlight something is code, but not run it, you can do
something like \texttt{mean(data\$column)}.

If you want to actually run R code, you need to modify it slightly to
insert code here. For instance, we could do ``the mean of the mpg column
is 20.090625

the standard deviation of the price column is 3989.4397381

\begin{Shaded}
\begin{Highlighting}[]
\NormalTok{correlationTest}\OtherTok{\textless{}{-}} \FunctionTok{cor.test}\NormalTok{(mtcars}\SpecialCharTok{$}\NormalTok{mpg, mtcars}\SpecialCharTok{$}\NormalTok{hp)}
\end{Highlighting}
\end{Shaded}

the correlation from \texttt{cor.test()} is -0.7761684

the correlation is -0.7761684

\subsubsection{Practice:}\label{practice-1}

Try Writing a sentence with any calculation you'd like in it.

\paragraph{More Resources}\label{more-resources-1}

More information about writing APA results can be found here

\subsection{Learning Outcome \#3:Create Write Ups that combine graphs,
interpretation, and analyses into a streamlined
document.}\label{learning-outcome-3create-write-ups-that-combine-graphs-interpretation-and-analyses-into-a-streamlined-document.}

Of course, we likely want to include tables and graphs in our write ups.
We can do this, but it is helpful to use multiple code chunks and save
information as needed.

\subsection{Example:}\label{example}

\paragraph{Results:}\label{results}

In the current study, we had 344 penguins participate. We collected
average flipper lengths, which can be seen in table 1.

\textbf{Table 1}

\begin{longtable}[]{@{}lr@{}}
\caption{average flipper length}\tabularnewline
\toprule\noalign{}
species & meanFlipper \\
\midrule\noalign{}
\endfirsthead
\toprule\noalign{}
species & meanFlipper \\
\midrule\noalign{}
\endhead
\bottomrule\noalign{}
\endlastfoot
Adelie & 189.9536 \\
Chinstrap & 195.8235 \\
Gentoo & 217.1870 \\
\end{longtable}

\textbf{bad table}

\begin{Shaded}
\begin{Highlighting}[]
\NormalTok{penguinFlipper}
\end{Highlighting}
\end{Shaded}

\begin{verbatim}
## # A tibble: 3 x 2
##   species   meanFlipper
##   <fct>           <dbl>
## 1 Adelie           190.
## 2 Chinstrap        196.
## 3 Gentoo           217.
\end{verbatim}

\subsubsection{Practice:}\label{practice-2}

As a class, let's come up with something we want to analyze and create
results for from the penguins dataset.

\begin{itemize}
\tightlist
\item
  What's the research question?
\item
  descriptive statistics?
\item
  inferential statistic?
\item
  how to present results? (table, graph, both?!)
\end{itemize}

\subsubsection{APA version of tables:}\label{apa-version-of-tables}

We don't always want just a table of counts, or we may want to format
statistical output in APA. To do so, the library \texttt{apaTables} has
code to report several different types of statistics. Code can be found
here

\end{document}
